\documentclass{article}

\begin{document}

\begin{center}
\textbf{審査会配布用の要旨}\\
\end{center}
題目:
画像処理と次数制約付き最小全域木に基づく細胞系統擬似時間の予測手法の改良\\
氏名:
岩城拓磨\\\\
要旨:\\
生物の細胞は多種多様な種類から構成されており、1つの細胞が分裂を繰り
返すことで多様性を持った細胞群へと分化していく。
このような細胞の多様性を生み出す運命決定の仕組みについては、いまだ多くの謎が残されている。
細胞の分化や発生のメカニズムを解明するには、各過程で個々の細胞がどのような振る舞いをしているかを正確に把握することが不可欠である。
シングルセルRNAシーケンシングは、個々の細胞レベルで遺伝子発現プロファイルを取得できる画期的な技術である。
シングルセルRNAシーケンシングを用いることで、個々の細胞における遺伝子発現データを大量に取得・解析することが可能となり、
細胞の分化プロセスを従来よりも詳細に追跡することが可能になった。
各細胞の内部状態の変化は時間的に同期されているわけではなく、各細胞ごとの時間軸で細胞は変化していく。
しかし研究者が欲しい情報は初期の細胞がどのような分化過程を辿るのかであるので、細胞全体として細胞の変化を追う必要がある。
それらを解析するために擬似時間解析は使用される。
擬似時間解析とは細胞の分化過程や疾患の進行などにおける細胞状態の遷移を理解するため、
遺伝子発現データに基づき、細胞間の類似性を指標として擬似的な時間軸に沿って並び替える手法である。
高次元データを擬似時間という1次元の特徴として捉えることであり、分化や疾患などの時系列の挙動を要約することができる。
擬似時間解析ツールのSlingshotはクラスタ済みのシングルセルRNAシーケンシングデータに対して最小全域木を構築した後PrincipalCurveで曲線にフィッティングすることで擬似時間を解析する。
しかしSlingshotは「細胞の状態が一度に3つ以上には遷移しない」という生物学的前提を考慮していない。
また入力データをクラスタリングする必要があるが、適切なクラスタ方法やパラメータの決定方法などは不明である。
本研究ではSlingshotのこれら課題について改良を加える。
クラスタ数を自動決定するために画像処理を用いる。
2次元に次元削減後のシングルセルRNAシーケンシングデータをヒストグラムを用いて画像として認識する。
その画像に対して細線化処理をすることで大まかな木構造を計算する。
その木構造の分岐と葉の数を調べることでシングルセルRNAシーケンシングデータのクラスタ数を決定しKmeansでクラスタを行う。
「細胞の状態が一度に3つ以上には遷移しない」という制約をSlingshotに組み込むために次数制限付き最小全域木を用いる。
最小全域木を構築する際にクラスタの各重心頂点の次数の数を３以下に制限し得られた最小全域木をprincipalcurveで曲線にフィッティングすることで細胞系統擬似時間の軌道とする。
この改良の有効性は、Bonemarrowmononuclearcell、Hematopoiesis、Humanfetalimmunecell、およびC.elegansの公開データセットを用いて実証した。
\end{document}
